%% 10_dsl.tex -- Chapter 10: Input Impedance, DC Servo Loops and Offset Rejection
%% Status: OUTLINE (build stage 3). Replace the outline with the full chapter when it is written.
\chapter{Input Impedance, DC Servo Loops and Offset Rejection}\label{ch:dsl}
\chapterlead{Electrodes bring hundreds of millivolts of DC and megaohms of source impedance. The front-end must reject the first without loading the second.}

\begin{specnote}[Chapter status: outline]
This chapter is an outline. It will be written in build stage 3 (Part II with Labs 3--5). The topics, the design question and the labs below are fixed; the derivations, worked examples and figures are still to come.
\end{specnote}

\begin{objectives}
\begin{itemize}
  \item derive the input impedance of a capacitively coupled, chopped front-end;
  \item design a positive-feedback impedance-boosting loop;
  \item design a DC servo loop (DSL) that rejects \qty{\pm 300}{\milli\volt} of electrode offset;
  \item compare DC-coupled and AC-coupled front-ends;
  \item state the trade-offs for dry electrodes.
\end{itemize}
\end{objectives}

\begin{insight}[Design question]
How do you tolerate \qty{\pm 300}{\milli\volt} of electrode offset and keep the input impedance high enough for dry electrodes?
\end{insight}

\section{Input impedance of chopped CCIAs}\label{sec:ch10-1}
Switched-capacitor input resistance, its frequency dependence, and the \qty{2.5}{\mega\ohm} the ECG standards require (\cref{ch:interference}).

\section{Positive-feedback impedance boosting}\label{sec:ch10-2}
Pre-charging the input capacitors from the output to cancel the current drawn from the electrode; stability and matching limits.

\section{The DC servo loop}\label{sec:ch10-3}
An integrator in a feedback path subtracts the electrode offset at the input; the loop sets the high-pass corner and must not add noise in band.

\section{DC-coupled versus AC-coupled front-ends}\label{sec:ch10-4}
DC coupling keeps the electrode offset for motion-artefact detection and needs a wide-range ADC; AC coupling removes it in the analog domain. The choice decides between the two converters of Part~III.

\section{Dry electrodes}\label{sec:ch10-5}
Megaohm source impedance, larger offsets and motion artefacts; what each front-end option buys and costs.

\begin{tryit}[Labs for this chapter]
Lab~5 (\cref{sec:lab5}) measures the input impedance with and without a chopper and verifies \qty{\pm 300}{\milli\volt} offset tolerance with the DC servo loop.
\end{tryit}

\paragraph{Read with this chapter} Yazicioglu, Van Hoof and Puers~\cite{yazicioglubook}, Spinelli et al.~\cite{spinelli2003}, Chi et al.~\cite{chi2010}.
